% Part: computability
% Chapter: recursive-functions
% Section: pr-functions-computable

\documentclass[../../../include/open-logic-section]{subfiles}

\begin{document}

\olfileid{cmp}{rec}{cmp}
\olsection{بنسټيزې بازګشتي تابعې محاسبه کېدونکې دي}

فرض کړئ يوه تابع~$h$ په بنسټيز بازګښت تعريف شوې ده:
\begin{eqnarray*}
h(\vec x, 0)     & = & f(\vec x) \\
h(\vec x, y + 1) & = & g(\vec x, y, h(\vec x, y))
\end{eqnarray*}
او فرض کړئ تابعې~$f$ او~$g$ محاسبه کېدونکې دي۔ (موږ~$\vec x$ د
$x_0$، \dots،~$x_{k-1}$ د لنډون دپاره کاروو۔) بيا
$h(\vec x, 0)$ په څرګند ډول محاسبه کېدے شي، ځکه دا يوازې
$f(\vec x)$ ده چې محاسبه کېدونکې مو فرض کړې ده۔ بيا
$h(\vec x, 1)$ هم محاسبه کېدے شي، ځکه~$1 = 0 + 1$ دے او په دې ډول
$h(\vec x, 1)$ يوازې دا ده:
\begin{align*}
 h(\vec x, 1) & = g(\vec x, 0, h(\vec x, 0)) =  g(\vec x, 0, f(\vec x)).
\intertext{په همدې ډول مخکښې تللے شو او دا محاسبه کولے شو:}
h(\vec x, 2) & = g(\vec x, 1, h(\vec x, 1)) = g(\vec x, 1, g(\vec x, 0, f(\vec x)))\\
h(\vec x, 3) & = g(\vec x, 2, h(\vec x, 2)) = g(\vec x, 2, g(\vec x, 1, g(\vec x, 0, f(\vec x))))\\
h(\vec x, 4) & = g(\vec x, 3, h(\vec x, 3)) = g(\vec x, 3, g(\vec x, 2, g(\vec x, 1, g(\vec x, 0, f(\vec x)))))\\
& \vdots
\end{align*}
نو په عمومي ډول د~$h(\vec x, y)$ د محاسبې دپاره په پرله‌پسې توګه
$h(\vec x, 0)$، $h(\vec x, 1)$، \dots محاسبه کوو، تر هغه چې
$h(\vec x, y)$ ته ورسېږو۔

نو که تابعې~$f$ او~$g$ محاسبه کېدونکې وي، يو بنسټيز بازګشتي تعريف
يوه نوې محاسبه کېدونکې تابع ورکوي۔ همدارنګه، که تابعې~$f$ او~$g_i$
محاسبه کېدونکې وي، د تابعو ترکيب هم محاسبه کېدونکې تابع ورکوي۔

څرنګه چې بنسټيزې تابعې~$\Zero$، $\Succ$ او~$\Proj{n}{i}$ محاسبه
کېدونکې دي، او ترکيب او بنسټيز بازګښت له محاسبه کېدونکو تابعو څخه
محاسبه کېدونکې تابعې ورکوي، نو هره بنسټيزه بازګشتي تابع محاسبه
کېدونکې ده۔

\end{document}
